\chapter*{怎麼用這本書}
\phantomsection\addcontentsline{toc}{chapter}{怎麼用這本書}
\markboth{怎麼用這本書}{}

這本書收錄研究所入學考試「資料結構與演算法」的原題，每一題都附完整詳解。 它是
《資料結構與演算法：C 語言、虛擬碼與研究所考試》（以下簡稱\textbf{資結書}）的習題本：
章號和章名跟資結書一樣，第~12~章的題目就是資結書第~12~章雜湊的內容。

\section*{題目怎麼排}
\phantomsection\addcontentsline{toc}{section}{題目怎麼排}

\begin{itemize}
\item \textbf{依主題}：一題歸到它主要考的那一章。 一題同時考兩章時，放在比較吃重的那一章。
\item \textbf{章裡依學校、學年度排}：同一份考卷的題目放在一起，照原卷的題號順序。
\item \textbf{想整份考卷計時練習}：附錄「依考卷查題」列出每一份考卷的每一題在第幾頁。
\item 合卷（資料結構和作業系統、計算機組織等科目印在同一份）只收資料結構與演算法的題目。
\end{itemize}

\section*{題目和詳解}
\phantomsection\addcontentsline{toc}{section}{題目和詳解}

\begin{itemize}
\item 題目照原卷：英文就是英文，記號、配分、選項照錄，錯字也不改。 原卷的圖用 TikZ 重畫。
\item 詳解是本書作者寫的，不是各校公布的答案。 選擇題逐個選項說明對錯，追蹤題寫出每一步，
證明題寫完整的證明，設計題寫虛擬碼、正確性理由和複雜度。 每個數值答案都用程式驗算過。
\item 原題有錯字、條件不足或有兩種讀法時，詳解第一段會說明，並寫出採用的解讀。
\item 詳解裡的「資結書第~N~章」「資結書第~N.M~節」指資結書的章節，虛擬碼的寫法也照資結書
第~1~章的約定。
\end{itemize}

\section*{建議的用法}
\phantomsection\addcontentsline{toc}{section}{建議的用法}

讀完資結書的一章，就做這本書的同一章。 先自己寫，寫完再看詳解；對答案時除了答案本身，也看
詳解的寫法（追蹤畫到哪一步、證明寫到多細），那是閱卷老師看的東西。 考前一兩個月，改用附錄
的「依考卷查題」整份計時做，練時間分配。

\section*{版權}
\phantomsection\addcontentsline{toc}{section}{版權}

依法令舉行之考試試題，依著作權法第~9~條不得為著作權之標的；本書收錄的題目來自各校公開的
研究所入學考試試題。 詳解與排版以 CC BY-NC-SA 4.0 授權。
